\section*{Chapitre \ref{ch:m2:distressed-sovereign-and-bank-capital-credit} --- Dette en difficulté, dette souveraine et crédit en fonds propres bancaires}

\begin{solution}{exo:m2:distressed-sovereign-and-bank-capital-credit:1}
Une décote sur le nominal de 40~\%. Les nouvelles obligations valent 47.92 à un taux de 8~\%, soit une
décote en valeur actualisée de 52.1~\%.
\end{solution}

\begin{solution}{exo:m2:distressed-sovereign-and-bank-capital-credit:2}
Une clause série par série exige une majorité qualifiée, généralement de 75~\%, de chaque
obligation prise séparément~; un \omterm{def:m2:distressed-sovereign-and-bank-capital-credit:cac}{créancier récalcitrant} (holdout) détenant une minorité de blocage dans une série peut donc maintenir cette série à l'écart.
Une clause agrégée à volet unique compte un seul vote sur l'ensemble des séries touchées, soit 75~\% de
leur principal total dans le modèle de l'ICMA, si bien qu'aucune série isolée ne peut bloquer, à condition que
toutes les séries se voient proposer les mêmes conditions.
\end{solution}

\begin{solution}{exo:m2:distressed-sovereign-and-bank-capital-credit:3}
Son déclencheur contractuel, généralement une baisse du ratio de fonds propres durs de la banque
sous un seuil fixé, qui dépend des pertes de la banque~; et le \omterm{def:m2:distressed-sovereign-and-bank-capital-credit:pon}{point de non-viabilité}, qui dépend
d'une décision de l'autorité jugeant que la banque ferait défaut sans dépréciation forcée ou
soutien public.
\end{solution}

\begin{solution}{exo:m2:distressed-sovereign-and-bank-capital-credit:4}
De 60.9~\% à 66.7~\%~: la clause lie les \omterm{def:m2:distressed-sovereign-and-bank-capital-credit:cac}{créanciers récalcitrants} aux deux tiers, si bien que la bande
se réduit de 8.3 points, à 5.8.
\end{solution}

\begin{solution}{exo:m2:distressed-sovereign-and-bank-capital-credit:5}
Cela signifiait que payer les créanciers ayant accepté l'échange imposait de régler
intégralement les \omterm{def:m2:distressed-sovereign-and-bank-capital-credit:cac}{créanciers récalcitrants} au même moment, et permettait à ces derniers de bloquer
les paiements par l'intermédiaire des banques qui les traitaient~; les créanciers acceptants perdaient la certitude d'être payés, ce qui
rendait plus difficile le placement des échanges futurs. Les contrats récents précisent que la clause
n'impose pas de paiement au prorata et recourent à des \omterm{def:m2:distressed-sovereign-and-bank-capital-credit:cac}{clauses d'action collective} agrégées.
\end{solution}

\begin{solution}{exo:m2:distressed-sovereign-and-bank-capital-credit:6}
La loi grecque ne pouvait modifier que les obligations de droit grec~; les obligations de droit anglais ne pouvaient pas être
modifiées sans le vote de leurs détenteurs, et un défaut sur celles-ci aurait exposé la Grèce
à des poursuites à l'étranger. Les \omterm{def:m2:distressed-sovereign-and-bank-capital-credit:cac}{créanciers récalcitrants} pesaient peu au regard de l'ensemble, si bien que les payer
intégralement revenait moins cher que le bras de fer~: le passager clandestin a été payé.
\end{solution}

\begin{solution}{exo:m2:distressed-sovereign-and-bank-capital-credit:7}
Refuser l'échange vaut alors 46.26 même si personne n'apporte ses titres, davantage que les 34.85 de
l'apport~: elle rapporte pour toute participation inférieure aux 75~\% de la clause. Une position juridique
plus forte élargit la bande à toutes les participations sous le seuil, si bien que seule la clause
empêche le passager clandestin d'en tirer profit.
\end{solution}

\begin{solution}{exo:m2:distressed-sovereign-and-bank-capital-credit:8}
Le classement s'applique en liquidation. Une obligation AT1 est aussi dépréciée ou convertie
selon ses propres conditions et par l'autorité au \omterm{def:m2:distressed-sovereign-and-bank-capital-credit:pon}{point de non-viabilité}, alors que la banque
est encore en activité, et rien dans ces conditions n'exige que les actionnaires soient effacés
au préalable~; dans le cas de Credit Suisse, les obligations AT1 ont été ramenées à zéro tandis que les actionnaires recevaient
3~milliards de francs suisses d'actions UBS. Les autorités européennes ont indiqué qu'elles appliqueraient
l'ordre de liquidation~; l'affaire suisse est devant les tribunaux.
\end{solution}

\begin{solution}{pb:m2:distressed-sovereign-and-bank-capital-credit:1}
\textbf{1.} 34.85 avec la part en numéraire, 29.85 sans.
\textbf{2.} 50~\% du nominal~; 65.2~\% en valeur actualisée.
\textbf{3.} 45.29 à 6~\%, 27.76 à 12~\%.
\textbf{4.} Parce que les nouvelles obligations versent un faible coupon pendant quinze ans~: leur valeur, et
donc la perte réelle des créanciers dépend de l'opinion du marché sur l'émetteur souverain après
l'échange.
\textbf{5.} 34.85 millions~: 29.85 millions en nouvelles obligations et 5 millions en numéraire.
\textbf{6.} Versé l'année suivante~: $100/1.09 = 91.74$~; le procès~: $130/1.09^6 = 77.51$~; en
défaut~: 15~; la branche sans paiement anticipé vaut $0.3 \times 77.51 + 0.7 \times
15 = 33.75$.
\textbf{7.} 33.75~; 34.73~; 38.97.
\textbf{8.} À 60.9~\%.
\textbf{9.} La clause le lie à l'échange~: il reçoit les nouvelles obligations, 29.85, sans
la part en numéraire.
\textbf{10.} 58.72.
\textbf{11.} Pas de façon stable~: pour toute participation dans cette bande, chaque créancier
préfère refuser l'échange, ce qui fait baisser la participation~; les créanciers qui pensent que 75~\% seront atteints apportent leurs titres pour
conserver la part en numéraire, ce qui la fait monter. Ce sont les anticipations qui tranchent.
\textbf{12.} La part en numéraire est la différence entre apporter ses titres et être lié, soit 5 pour 100~:
c'est ce que perd le \omterm{def:m2:distressed-sovereign-and-bank-capital-credit:cac}{créancier récalcitrant} en attendant, une fois la clause contraignante.
\textbf{13.} En achetant plus de 25~\% du principal agrégé, ce qui représente bien plus que
un quart d'une seule série~; l'agrégation rend une minorité de blocage bien plus coûteuse.
\textbf{14.} Ils étaient payés intégralement à mesure que leurs obligations arrivaient à échéance, par exemple 435 millions d'euros en
mai 2012, parce que le droit grec ne pouvait pas les atteindre et qu'ils étaient peu nombreux.
\textbf{15.} En juillet 2014, ils n'ont pas été payés, car les tribunaux interdisaient le paiement des
obligations échangées à moins que les \omterm{def:m2:distressed-sovereign-and-bank-capital-credit:cac}{créanciers récalcitrants} ne soient payés intégralement au même moment.
\textbf{16.} Cela dépend de la trésorerie de l'émetteur souverain, de la situation politique et de son exposition
juridique, ainsi que du nombre de créanciers qui refusent l'échange, autant d'éléments qui évoluent au cours de la restructuration.
\textbf{17.} Question d'appréciation~: les \omterm{def:m2:distressed-sovereign-and-bank-capital-credit:cac}{créanciers récalcitrants} font respecter les contrats et discipliner les débiteurs,
mais ils peuvent bloquer un allègement que la plupart des créanciers acceptent~; les clauses agrégées
laissent une large majorité décider.
\textbf{18.} Peu de chose est liquide~: les CDS sur l'émetteur souverain, s'il en existe qui ont été
dénoués sur un défaut, n'existent plus~; le fonds peut couvrir le rendement de sortie grâce aux autres obligations négociées du souverain ou
aux devises liées, mais l'issue juridique est impossible à couvrir.
\textbf{19.} Résultat retenu~: \emph{bande du passager clandestin}~: refuser l'échange rapporte pour une participation
comprise entre 60.9~\% et 75~\%~; à 75~\%, la clause lie le \omterm{def:m2:distressed-sovereign-and-bank-capital-credit:cac}{créancier récalcitrant} à l'offre, qui vaut 29.85
contre 34.85 pour l'apport.
\textbf{20.} Elles empêchent une minorité de profiter sans contrepartie d'une restructuration
qu'une majorité qualifiée a acceptée.
\end{solution}

\begin{solution}{iq:m2:distressed-sovereign-and-bank-capital-credit:1}
Une entreprise peut déposer son bilan, procédure où un tribunal impose un plan aux catégories de
créanciers et peut liquider des actifs~; un souverain ne le peut pas, dispose de peu d'actifs saisissables, et
se restructure par \omterm{def:m2:distressed-sovereign-and-bank-capital-credit:distressed}{offres d'échange} et par les termes des contrats. Les \omterm{def:m2:distressed-sovereign-and-bank-capital-credit:cac}{créanciers récalcitrants} peuvent
poursuivre en justice mais peinent à recouvrer, et l'issue dépend du droit applicable, des \omterm{def:m2:distressed-sovereign-and-bank-capital-credit:cac}{clauses d'action collective} et
de la politique.
\iqlookfor{pas de tribunal, contrats, \omterm{def:m2:distressed-sovereign-and-bank-capital-credit:cac}{créanciers récalcitrants}.}
\end{solution}

\begin{solution}{iq:m2:distressed-sovereign-and-bank-capital-credit:2}
Une obligation bancaire perpétuelle comptée en fonds propres de continuité d'exploitation, à coupons discrétionnaires,
dépréciée ou convertie lorsqu'un déclencheur de ratio de fonds propres est atteint ou lorsque l'autorité
constate le \omterm{def:m2:distressed-sovereign-and-bank-capital-credit:pon}{point de non-viabilité}. Risques~: suppression du coupon, absence de rachat
à la première date de call, dépréciation ou conversion, et choix de l'ordre des pertes
par les autorités, comme l'a montré 2023.
\iqlookfor{coupons, prorogation, déclencheurs, pouvoir discrétionnaire de l'autorité.}
\end{solution}

\begin{solution}{iq:m2:distressed-sovereign-and-bank-capital-credit:3}
Évaluer les nouveaux instruments à un rendement de sortie estimé à partir de crédits restructurés comparables,
ajouter les espèces et les éventuelles incitations, et comparer avec le recouvrement obtenu en refusant l'échange,
un arbre de probabilités portant sur le fait d'être remboursé, de poursuivre en justice ou de rester en défaut, actualisé
de façon cohérente~; le choix dépend de la participation attendue et des modalités juridiques.
\iqlookfor{rendement de sortie, alternative du refus, participation.}
\end{solution}

\begin{solution}{iq:m2:distressed-sovereign-and-bank-capital-credit:4}
Chaque créancier préfère que les autres acceptent la perte tout en conservant l'intégralité de sa créance, si bien que
la participation peut stagner. Les \omterm{def:m2:distressed-sovereign-and-bank-capital-credit:cac}{clauses d'action collective} permettent à une majorité qualifiée d'imposer sa décision
à la minorité~; les clauses agrégées comptent l'ensemble des obligations, de sorte qu'une minorité de blocage dans une
seule obligation ne suffit plus.
\iqlookfor{comportement de passager clandestin et vote contraignant.}
\end{solution}

\begin{solution}{iq:m2:distressed-sovereign-and-bank-capital-credit:5}
Que les clauses contractuelles et le droit d'urgence, et pas seulement l'ordre de liquidation, déterminent qui
supporte la perte~; que les autorités peuvent différer d'un cas à l'autre, comme l'a montré la déclaration
de l'Union européenne~; et que le fondement juridique peut être contesté pendant des années. Les investisseurs peuvent
désormais exiger une prime au titre de ce risque juridictionnel.
\iqlookfor{contrat face à hiérarchie, juridiction, risque juridique.}
\end{solution}

\begin{solution}{iq:m2:distressed-sovereign-and-bank-capital-credit:6}
Extraire du prospectus de chaque obligation un enregistrement comprenant~: le droit applicable, la formulation de la \omterm{def:m2:distressed-sovereign-and-bank-capital-credit:pp}{clause
pari passu}, le type de clause (série par série, à deux volets, à volet unique), les seuils, le quorum et les titres privés
de droit de vote~; tenir à jour les montants en circulation à partir des données d'émission et de rachat~; calculer les
votes nécessaires selon chaque procédure ainsi que la participation qui bloquerait chaque série~;
signaler les modifications lorsque de nouvelles obligations sont émises.
\iqlookfor{données contractuelles structurées et arithmétique des votes.}
\end{solution}
